\documentclass[11pt,letterpaper]{article}
\usepackage{hanzo-defense}

\doctitle{Agentic Engineering at Regulated Scale:\\Building a Post-Quantum, FHE-Native\\Securities Platform}
\docid{HAI-DEF-2026-013}
\docversion{Version 1.0}
\docdate{2026}
\docauthors{Zach Kelling --- Hanzo Team}
\docorgs{Hanzo Industries\\research@hanzo.ai}

\begin{document}
\makecover

\begin{abstract}
We document the engineering capability behind Hanzo's stack by way of its hardest
production proof: an agent-driven team built a fully regulated securities
venue---an SEC-registered alternative trading system (ATS) with broker-dealer and
transfer-agent functions---on a next-generation foundation of production
post-quantum cryptography, a high-performance native-GPU blockchain, and
fully-homomorphic-encryption (FHE) privacy. The build \emph{consolidated} rather
than sprawled---200+ microservices collapsed into 3 compiled binaries, with 3.07M
lines of code reduced to 1.06M active ($\sim$3.5M lines of dead code deleted)---reaching
production in record time with a small team, and passed FINRA/SEC approval and SOC\,2
compliance. This is not template web
software; it is among the most compliance-heavy and cryptographically advanced
software domains that exists. We describe the throughput of the agentic
engineering method that produced it, and we expand on the defense relevance of
its FHE layer---computing on data one is not permitted to see.
\end{abstract}

\paragraph{Executive Summary.}
The fastest way to judge an engineering organization is what it has shipped under
real constraints. Hanzo's founding CTO, Zach Kelling, was ranked the
\textbf{\#1 user of Anthropic by token count worldwide in 2025} (per public usage
tracking at \texttt{claudecount.com})---a marker of an engineer operating
agentic AI at the frontier of throughput. That throughput is not spent on
brochure sites: the flagship production system built with Hanzo's agentic stack
is a \textbf{fully regulated securities exchange}---an SEC-registered ATS with
broker-dealer and transfer-agent functions---running on \textbf{post-quantum
cryptography}, a \textbf{native-GPU high-performance blockchain}, and
\textbf{fully homomorphic encryption} for privacy. The build \emph{consolidated}
rather than sprawled---200+ microservices collapsed into 3 compiled binaries, 3.07M
lines of code cut to 1.06M active---shipped fast with a small team, and cleared
FINRA/SEC approval and SOC\,2. Regulated financial infrastructure is one of the hardest,
most audited software domains there is; clearing it on a cryptographically
advanced foundation is the proof that agent-driven engineering is production-grade,
not a demo---and that it is exactly the capability needed to modernize defense
software stacks.

\tableofcontents
\newpage

% ============================================================================
\section{The Engineer and the Throughput}
% ============================================================================

Agentic engineering multiplies a strong engineer; it does not replace one. Hanzo's
founding CTO works at the frontier of both AI/ML and cryptography---the two
disciplines this stack fuses---and operates agentic AI at a volume that placed him,
by public token-count tracking (\texttt{claudecount.com}), as the top Anthropic
user in the world in 2025. The relevant point is not the ranking for its own sake;
it is what that volume was spent producing. Raw AI usage that yields disposable
prototypes proves nothing. Raw AI usage that yields a regulated, audited,
cryptographically advanced production system proves the method works where it is
hardest.

% ============================================================================
\section{What Was Built}
% ============================================================================

The flagship system is a regulated securities venue, not a content site. In U.S.
securities terms it spans three regulated functions that ordinarily take separate
firms years to stand up:

\begin{itemize}
  \item \textbf{Alternative Trading System (ATS)} --- an SEC-registered trading
  venue (Regulation ATS), matching and executing orders.
  \item \textbf{Broker-Dealer (BD)} --- FINRA-regulated, with the supervisory,
  recordkeeping, and capital obligations that entails.
  \item \textbf{Transfer Agent (TA)} --- SEC-registered, maintaining the
  authoritative record of securities ownership and transfers.
\end{itemize}

These functions are governed by some of the most exacting compliance regimes in
software. The notable engineering fact is that the build \emph{consolidated} rather
than sprawled: a legacy estate of \textbf{200+ microservices} on a public cloud was
collapsed into \textbf{3 compiled binaries}, and the codebase went from
\textbf{3.07M lines to 1.06M active} ($\sim$70\% smaller, with $\sim$3.5M lines of
dead code deleted; the trading/auth core itself shrank from $\sim$264{,}000 to
$\sim$44{,}000 lines, with 1.3\,GB of vendored dependencies eliminated and a single
48\,MB binary subsuming 17+ formerly separate services). It reached production
\textbf{in record time}, with a \textbf{small team}, and was carried through
\textbf{FINRA/SEC approval} and \textbf{SOC\,2} compliance---leaner than the system
it replaced, not larger.

% ============================================================================
\section{One Engineer, a Team's Output}
% ============================================================================

The result that matters most is the ratio of output to headcount. The agentic method
let a \emph{single} AI-augmented engineer carry the bulk of a regulated production
platform that a conventional organization staffs as a full team---and the platform's
own engineering audit documents it. Across an \textbf{81-day} window spanning
\textbf{334 repositories} and \textbf{31{,}058 commits}, the audit attributed
\textbf{15{,}294 commits} and \textbf{114M raw git-log lines}---framework-wide refactors,
dependency upgrades, automation, and production code---to the agentic CTO workflow,
measured at \textbf{\$0.0006 per line} by the audit's methodology. Within that, a focused
\textbf{10-day remediation sprint} produced \textbf{3{,}694 commits}, removed
$\sim$\textbf{3.5 million lines of dead code} across 181 repositories, and authored
\textbf{100\% of the critical- and high-severity security fixes}---single-handedly---while
consolidating 200+ microservices into three binaries. The security and compliance gaps
were closed on the path to approval (a source audit of \textbf{93 findings}, 22 critical,
plus \textbf{780+ dependency vulnerabilities} remediated to zero actionable), through
FINRA/SEC approval and SOC\,2.

The economic consequence is direct and measured. When one engineer, multiplied by the
agentic stack, produces what a team produces, the \textbf{cost per line of shipped,
audited, production code} collapses---here, \textbf{\$0.0006 per line} by the audit's own
methodology, orders of magnitude below conventional team rates. For a defense or government program the implication is the whole argument:
the same regulated, post-quantum, FHE-native system, delivered by a fraction of the
headcount at a fraction of the cost---and owned, not rented. This is what ``agentic
engineering'' means in practice---not a chatbot writing snippets, but one expert
operating at the throughput of a department.

% ============================================================================
\section{The Foundation: Post-Quantum, GPU Blockchain, FHE}
% ============================================================================

What makes the build a genuine engineering flex is the foundation it runs on---
each element of which is independently hard:

\begin{itemize}
  \item \textbf{Production post-quantum cryptography.} All three NIST PQC
  standards (ML-KEM/203, ML-DSA/204, SLH-DSA/205) in a live system with hybrid
  classical/quantum security---so the chain of custody for securities records is
  protected against a future quantum adversary.
  \item \textbf{Native-GPU high-performance blockchain.} A chain serving as the
  authoritative, immutable, post-quantum-anchored record of transfers, with
  GPU-accelerated cryptographic operations for throughput.
  \item \textbf{Fully homomorphic encryption (FHE).} GPU-accelerated FHE
  (BFV/BGV/CKKS/TFHE) with threshold decryption, so privacy-sensitive computation
  runs \emph{on ciphertext}---no single component ever holds both the plaintext
  and the result.
\end{itemize}

A team that can compose post-quantum cryptography, a GPU blockchain, and FHE into
a system that clears securities regulators is a team that can modernize a defense
software stack to the same foundation.

% ============================================================================
\section{FHE for the Military: Computing on Data You Cannot See}
% ============================================================================

Fully homomorphic encryption is the capability with the sharpest defense upside,
so it is worth drawing out. FHE allows arbitrary computation directly on encrypted
data: the operator performs the computation and returns an encrypted result
without ever decrypting the inputs, and only the data owner (or a threshold of
keyholders) can decrypt the output. The system doing the work never sees the
plaintext. In the securities platform this protects order and ownership data; the
same primitive answers several defense problems that have no clean classical
solution:

\begin{itemize}
  \item \textbf{Cross-domain transfer on ciphertext.} A cross-domain guard can
  evaluate a release policy against encrypted content and decide whether data may
  move between security levels \emph{without ever seeing the data}---removing the
  guard itself as a point of compromise (developed in the cross-domain paper).
  \item \textbf{Coalition analytics without disclosure.} Allies can compute a
  joint result---shared targets, overlapping indicators, combined logistics---over
  their combined data while each party's raw data stays encrypted and private.
  Nobody has to surrender sources or methods to get the joint answer.
  \item \textbf{Private intelligence fusion.} Multi-source correlation can run
  while the sources remain encrypted, so a breach of the fusion system does not
  expose the underlying collection.
  \item \textbf{Encrypted ML inference.} A model can classify encrypted imagery or
  signals: the data owner never exposes the input, and the model operator never
  sees it---useful when the data is sensitive, the model is sensitive, or both.
  \item \textbf{Computation on untrusted infrastructure.} Sensitive workloads can
  run on contractor or shared infrastructure without exposing the data to the host.
\end{itemize}

FHE's historical objection is speed; the practical answer in this stack is
GPU acceleration of the underlying number-theoretic transforms, plus threshold
decryption so no single party holds the key. ``Compute on data you are not allowed
to see'' is a capability defense has wanted for decades; here it is a production
component, not a research slide.

% ============================================================================
\section{Why Open Source Was the Accelerant}
% ============================================================================

The obvious question is how a small team shipped a regulated, cryptographically
advanced system of this scale so quickly. The answer is open source---specifically,
an \emph{owned} open-source stack composed by an agentic engineering method.

Hanzo did not assemble the platform from a patchwork of third-party SaaS. It
composed it from coherent open-source building blocks that Hanzo itself authors
and controls---identity and access management, key management, the data layer, the
gateway and ingress, the ZAP transport, the post-quantum and FHE cryptography, the
operator, and the Zen models---each reusable, inspectable, and owned. The agentic
stack is the multiplier; the owned OSS foundation is the leverage. You do not
rebuild identity, key management, or a gateway for each new system; you compose
owned, hardened components and point the agents at the application on top. That is
how a regulated platform of this scope reaches production fast---and lean.

For defense, open source inverts the usual security argument. A closed system asks
you to trust a vendor you cannot inspect; an open one lets you read every line,
verify it---here, with machine-checked proofs---confirm there are no hidden
backdoors, and run it on hardware you control. Auditable beats opaque whenever the
stakes are high. Open source is therefore not a cost-cutting compromise: it is
simultaneously the source of \textbf{velocity} (compose, do not rebuild), of
\textbf{security} (inspectable and formally verifiable), and of \textbf{sovereignty}
(own it, modify it, run it anywhere, including air-gapped). The three properties a
defense program most wants are the same three that made the build fast.

Stated plainly as a hook: \emph{Hanzo builds massive, secure systems quickly
because it builds on an open-source stack it owns---so the speed, the security, and
the sovereignty all come from the same source.}

% ============================================================================
\section{Why It Matters for Defense Modernization}
% ============================================================================

The throughput and the difficulty compound. The same agentic engineering method
that shipped a regulated, post-quantum, FHE-native securities exchange fast and
small is the method Hanzo brings to defense and government modernization: take an
existing, already-accredited stack and modernize it in place---to a sovereign,
post-quantum, self-hostable footprint with FHE-grade privacy---faster and cheaper
than a traditional rebuild, and owned by the operator. The securities venue is the
existence proof that the method holds up under the most demanding regulatory and
cryptographic constraints civilian software offers; defense modernization is the
same engine pointed at the mission.

\end{document}
